\section{The kernel}
\label{sec:kernel}

\subsection{The matrix--vector product}

The kernel keeps the geometry of our earlier one: one warp per output row, a
256-thread CTA holding 8 rows, the activation staged in shared memory in tiles
of $T$ blocks ($96T$ bytes) with two barriers per tile, one lane per block, a
shuffle reduction, and an epilogue that adds the exactly-kept tail and
multiplies by $\sigma_{\mathrm{row}}$. Only the decode changed.

Three choices make a six-load decode competitive with a shift-and-mask stream.
First, table values are stored as biased bytes, $\mathrm{val}+128$. Placed
under \texttt{0x4b0000\_\_}, such a byte reads as the float
$2^{23}+\mathrm{byte}$ exactly, and one subtraction recovers the value, so a
coordinate costs one \texttt{prmt}, one \texttt{FADD} and one \texttt{FFMA}
with no integer-to-float conversion. Second, every loop is unrolled with
constant indices, so nothing spills: the compiler reports 40 registers and no
local memory, the same as \Planes{}. Third, the only data-dependent reads are
the six table loads, so lanes never diverge. Counted from the source, a block
costs about 141 instructions, of which 54 on the FMA pipe and 85 on the ALU
pipe, against about 380 plus 24 conversions for a decoder that extracts each
coordinate and converts it.

Because the kernel reads the file's own words there is no transcoding pass at
load. The served object loads in 12.1~s where the f16 checkpoint loads in
81.8~s in the same process.

\subsection{Where the time goes when the tile changes}
\label{sec:tile}

On an SM, L1 and shared memory come out of one block of SRAM: on Ada (sm\_89),
128~KB per SM, of which at most 100~KB can be assigned to shared
memory~\citep{cudaguide}. The activation tile is shared memory. The decoder
tables are not; they are ordinary global reads, and what shared memory takes,
L1 does not get. So the tile size, which changes no bit of output and costs no
bit of storage, changes how much L1 is left for an 18.4~KiB table.

Occupancy does not change with it. Both arms launch 256 threads at 40
registers, so six CTAs are resident at every tile
($6\times256\times40 = 61{,}440$ of 65{,}536 registers), and shared memory is
never the binding limit ($6\times12{,}288 < 102{,}400$~B). What the tile does
change, besides the carve-out, is the loop structure: barriers per row and trip
counts. That part is common to both arms.

\begin{figure}[t]
\centering
\includegraphics[width=0.60\linewidth]{fig_tile.pdf}
\Description{Median milliseconds against tile size 32, 64 and 128 for three
arms. Planes14 is flat near 5 ms across the three tiles, Tetra falls from 4.08
at tile 128 to 3.42 at 64 and rises slightly at 32, and the no-weights control
is flat near 2.3 ms.}
\caption{The tile sweep on the served path, one process, one L40S. Each curve
is one arm against itself across tiles; the percentage is that arm's own
amplitude, largest over smallest. Data: tuile-l40s.csv, from
tuile-l40s-2026-09-20.}
\label{fig:tile}
\end{figure}

\begin{table}[t]
\centering\footnotesize\setlength{\tabcolsep}{5pt}
\begin{tabular}{rrrrr}
\toprule
$T$ & shared per SM & \texttt{nullk} & \Planes{} & \Tetra{}\\
\midrule
128 & 72~KB & 2.272 & 5.070 & 4.078\\
64 & 36~KB & 2.288 & 4.995 & \textbf{3.423}\\
32 & 18~KB & 2.388 & 5.044 & 3.553\\
\midrule
\multicolumn{2}{l}{amplitude} & 5.1\,\% & 1.5\,\% & 19.1\,\%\\
\bottomrule
\end{tabular}
\caption{The numbers behind Fig.~\ref{fig:tile}: median ms over 252
projections, one process, one L40S. \texttt{nullk} keeps the launch geometry
and reads no weight byte. Shared memory per SM is six resident CTAs times
$96T$ bytes. Data: tuile-l40s.csv.}
\label{tab:tile}
\end{table}

Across $T \in \{128, 64, 32\}$ the \Planes{} arm moves 1.5\,\%, the
no-weights control 5.1\,\%, and the \Tetra{} arm 19.1\,\%
(Fig.~\ref{fig:tile}, Table~\ref{tab:tile}). The three arms share the launch
geometry, the staging code and the occupancy; the \Tetra{} arm alone reads
tables. We read the residue as pressure on the residency of those tables in
L1, and we act on it: the served default is now $T=64$ on sm\_89. We have not
demonstrated it. The rented platform denies performance counters
(\texttt{ERR\_NVGPUCTRPERM}), so no cache hit rate or per-level traffic was
measured, and a control that shares the schedule without reading tables cannot
rule out an interaction between the two. On sm\_120 the same sweep prefers
$T=32$, so the best tile is a property of the card rather than of the format.
All kernel numbers in \S\ref{sec:results} are at $T=64$.
